\section{Ethics statement}\label{app:ethics}

\pseudopara{Risks associated with this paper.}
This paper's contribution can be divided into three parts, each with its own risks: a new dataset of forbidden prompts to use in jailbreak evaluation, a new automated evaluator to evaluate victim model responses, and an experimental analysis of previously published jailbreaks from the literature.

In discussing these risks, it is worth noting three things.
First, the jailbreak techniques that we discuss in this paper are aimed at making ``aligned'' models give advice on topics that they were trained not to give advice on.
This is only one type of vulnerability, and the attacks that we use are not necessarily useful for other tasks like indirect prompt injection of LLM-based applications ~\citep{greshake2023not}.
Second, unaligned open-source models like Dolphin ~\citep{dolphin} are already freely available but are less capable than leading closed-source models.
Thus, the misuse potential of prompt-based jailbreaks mostly comes from the possibility of exploiting deployed closed-source models.
Third, to the best of our knowledge, the real-world harm caused by jailbreaks has so far been limited to minor reputational damage to technology companies.
As a result, broader negative social impacts of jailbreaking---like use in terrorism~\cite{westpoint2024generating}---remain hypothetical for now.

Given these considerations, we believe that the most important risks associated with releasing the components of this paper are as follows:

\begin{itemize}
  \item \textbf{StrongREJECT forbidden prompts.} Our dataset consists of both novel and existing forbidden prompts. Many of these prompts contain premises that are offensive or implicitly suggest possible ways to hurt others. We believe that the possible negative impact associated with the release of the forbidden prompts is low since all of the prompts are either already in the public domain or were manually created by us to reflect the sort of harmful material that could be found on the web by a lay person with under an hour of searching. For example, many of our misinformation questions are inspired by real fake news articles.

  \item \textbf{StrongREJECT automated evaluator.} Releasing our automated evaluator provides both jailbreak researchers and malicious attackers with an improved way to gauge whether jailbreaks are effective.
This could be abused by, e.g., using the automated evaluator score as the maximization criterion for a jailbreak search algorithm like PAIR~\citep{chao2023jailbreaking}.
We have not tested whether our automated evaluator score is robust to optimization, so it is difficult to assess this risk, although in the worst case, it could make jailbreak search algorithms more effective for attackers, relative to using existing automated evaluators.

  \item \textbf{Our experimental analysis.} Our experimental analysis focuses on the relative performance of various published jailbreak techniques.
As with the automated evaluator, this information could be misused to better target malicious attacks on real-world LLMs.
However, the magnitude of potential risk remains low, since all of these techniques were already publicly available.
\end{itemize}

Since the expected damage of these three risks is small, we believe that they are outweighed by the positive impact of giving researchers an improved evaluation for jailbreaking.

\pseudopara{Data ethics.} We sourced questions from publicly available datasets and refer users to the original datasets' licenses in our codebase. Custom data was generated by the authors.